\chapter{Introduction}
\label{chap:intro}

\section{Background and Motivation}

Agriculture is heavily influenced by environmental stochasticity, resource limitations, and complex operational constraints. Deciding when to fertilize fields and how to plan crop rotations over long horizons are central challenges for maximizing agricultural yield while minimizing ecological impact and operational costs. Deep Reinforcement Learning (DRL) has emerged as a promising method to solve such sequential decision-making problems, allowing agents to learn optimal policies through repeated interactions with simulated environments. This report outlines the final capstone experiments of the thesis, assessing how different DRL algorithms respond to various environmental setups, shaping functions, and exploration hyperparameters in agricultural simulations.

\section{Scope of the Report}

This comprehensive document synthesizes the two canonical experiment collections generated as the definitive final reporting pack of the thesis:
\begin{enumerate}
    \item \textbf{The Core 113-Run Study}: A large-scale matrix consisting of 113 formally curated runs covering crop planning, fertilization, grouped weather distributions, and guarding strategies for hierarchical RL.
    \item \textbf{The 42-Run Ablation Study}: An extensive, low-hanging hyperparameter ablation sweep probing three specific hypotheses on policy entropy, hierarchical shaping, and nutrient cost weighting.
\end{enumerate}

By exhaustively logging evaluation metrics, environment dynamics (e.g., episode lengths, normalized metrics), and step-by-step diagnostic curves, these experiments robustly establish the performance hierarchy between \textit{Advantage Actor-Critic} (A2C), \textit{Proximal Policy Optimization} (PPO), and \textit{Deep Q-Network} (DQN) baselines in deterministic vs. stochastic weather conditions.

\section{Research Objectives}

The primary objectives of the experiments detailed herein were to:
\begin{itemize}
    \item Establish a robust baseline comparison between on-policy algorithms (PPO, A2C) inside deeply uncertain agricultural environments.
    \item Analyze the statistical variance introduced by environment stochasticity (e.g., Fixed Weather vs. Random Weather) on convergence and policy stability.
    \item Explore specific ablation parameters: Policy Entropy ($\mathcal{H}$), Blocked Nutrient Compliance Shaping Penalties, and variations in Nutrient Cost Weights.
    \item Determine the safest, most robust default algorithm setups for real-world deployment or extended horizon simulations (e.g., $T=3000, 5000$ years).
\end{itemize}

Through rigorous statistical validation, pairing multiple seeds ($N=3$), this report aims to construct an absolute hierarchy of reinforcement learning approaches customized for the given agricultural domains.
